

\medskip
\noindent\SGAInlineRefTarget{sga3:iii:proposition:0-11}\textbf{Proposition 0.11.}
\textit{Let \(E\) be an \(S\)-scheme, and put
\(X=\underline{\operatorname{Aut}}_S(E)\). The kernel of the morphism of
\(S\)-group functors}
\[
 p_X:X\longrightarrow X^+
\]
\textit{is canonically identified with the commutative \(S\)-group functor
\(L'_X\) defined by}
\[
 \operatorname{Hom}_S(Y,L'_X)
 =\operatorname{Hom}_{\mathscr O_{E_0\times_{S_0}Y_0}}
 \left(
   \Omega^1_{E_0/S_0}\otimes_{\mathscr O_{S_0}}\mathscr O_{Y_0},
   \mathscr J\mathscr O_{E\times_S Y}
 \right),
\]
\textit{where \(Y\) denotes a variable \(S\)-scheme.}

Indeed, for a variable \(S\)-scheme \(Y\), we have
\[
 \operatorname{Hom}_S(Y,X)=\operatorname{Aut}_Y(E\times_S Y)
\]
and
\begin{align*}
 \operatorname{Hom}_S(Y,X^+)
  &=\operatorname{Hom}_S(Y_{\mathscr J},X)\\
  &=\operatorname{Aut}_{Y_{\mathscr J}}(E\times_S Y_{\mathscr J})\\
  &=\operatorname{Aut}_{Y_{\mathscr J}}
    \bigl((E\times_S Y)\times_Y Y_{\mathscr J}\bigr).
\end{align*}
Applying \SGARefLink{sga3:iii:corollary:0-7}{0.7}~\SGARefLink{sga3:iii:corollary:0-7:item:b}{(b)} to the \(Y\)-scheme \(E\times_S Y\), we obtain an
isomorphism of groups
\[
 \operatorname{Hom}_S(Y,L'_X)
 \simeq
 \operatorname{Ker}\!\left(
   \operatorname{Hom}_S(Y,X)\longrightarrow
   \operatorname{Hom}_S(Y,X^+)
 \right).
\]
This isomorphism is readily checked to be functorial in the \(S\)-scheme
\(Y\). Hence it yields an isomorphism of \(S\)-groups
\[
 L'_X\simeq\operatorname{Ker}(X\longrightarrow X^+),
\]
which completes the proof of \SGARefLink{sga3:iii:proposition:0-11}{Proposition 0.11}.

\medskip
\noindent\SGAInlineRefTarget{sga3:iii:corollary:0-12}\textbf{Corollary 0.12.}
\textit{Retain the notation of \SGARefLink{sga3:iii:proposition:0-11}{0.11}: \(E\) is an \(S\)-scheme and
\(X=\underline{\operatorname{Aut}}_S(E)\). There is a natural action \(f\)
of \(X\) on \(L'_X\), defined as follows. For every \(S\)-scheme \(Y\),}
\[
 \operatorname{Hom}_S(Y,X)=\operatorname{Aut}_Y(E\times_S Y)
\]
\textit{and}
\[
 \operatorname{Hom}_S(Y,L'_X)
 =\operatorname{Hom}_{\mathscr O_{E_0\times_{S_0}Y_0}}
 \left(
   \Omega^1_{E_0\times_{S_0}Y_0/Y_0},
   \mathscr J\mathscr O_{E\times_S Y}
 \right).
\]
\textit{Here}
\[
 \Omega^1_{E_0/S_0}\otimes_{\mathscr O_{S_0}}\mathscr O_{Y_0}
 \simeq \Omega^1_{E_0\times_{S_0}Y_0/Y_0}
\]
\textit{(see EGA~IV\textsubscript{4}, 16.4.5). The first group acts on
the second by transport of structure, and this action is functorial in
\(Y\). One then has}
\begin{equation}
 x\,i(m)\,x^{-1}=i\bigl(f(x)m\bigr),
 \tag{0.12.1}\label{sga3:iii:corollary:0-12:equation:0-12-1}
\end{equation}
\textit{for every \(Y\to S\), \(x\in\operatorname{Hom}_S(Y,X)\), and
\(m\in\operatorname{Hom}_S(Y,L'_X)\).}

Indeed, this follows from \SGARefLink{sga3:iii:corollary:0-7}{0.7}~\SGARefLink{sga3:iii:corollary:0-7:item:c}{(c)}, applied to the \(Y\)-scheme
\(E\times_S Y\).

\medskip
\noindent\SGAInlineRefTarget{sga3:iii:paragraph:0-13}\textbf{Reminder 0.13.}
The direct image of a quasi-coherent module under a morphism of finite
presentation is quasi-coherent. Under the same hypotheses, formation of
the direct image commutes with flat base change. In the situation

\SGAThreeNativeDiagram[0.50\textwidth]
  {native_diagrams/exp3/diagram_026_exp3_p020_reminder013_flat_base_change_square.tex}
  {The flat base-change square in Reminder \SGARefLink{sga3:iii:paragraph:0-13}{0.13}.}
  {fig:exp3-reminder-013-flat-base-change}
  {LNM 151, pp.~97--98}

if \(f\) (and hence \(f'\)) is of finite presentation and \(g\) (and
hence \(g'\)) is flat, then for every quasi-coherent
\(\mathscr O_T\)-module \(\mathscr F\),
\[
 f_*(\mathscr F)\otimes_{\mathscr O_S}\mathscr O_{S'}
 =f'_*\bigl(\mathscr F\otimes_{\mathscr O_S}\mathscr O_{S'}\bigr),
\]
or, more elegantly,
\[
 g^*\bigl(f_*\mathscr F\bigr)=f'_*\bigl(g'^*\mathscr F\bigr).
\]
These two facts hold more generally for a quasi-compact and
quasi-separated morphism \(f\); see EGA~I, 9.2.1, and
EGA~III\textsubscript{1}, 1.4.15 in the quasi-compact separated case
(taking EGA~III\textsubscript{2}, Err\textsubscript{III} 25 into
account), as well as EGA~IV\textsubscript{1}, 1.7.4 and 1.7.21.

\medskip
\noindent\SGAInlineRefTarget{sga3:iii:remark:0-14}\textbf{Remark 0.14.}\footnote{Editor's note: The first part of
this remark has been expanded, and a second part has been added.}
Return to the notation of \SGARefLink{sga3:iii:proposition:0-11}{0.11}: let \(E\) be an \(S\)-scheme, put
\(X=\underline{\operatorname{Aut}}_S(E)\), and let \(L'_X\) be the
commutative \(S\)-group functor defined by
\begin{align*}
 L'_X(Y)
 &=\operatorname{Hom}_{\mathscr O_{E_0\times_{S_0}Y_0}}
   \left(
     \Omega^1_{E_0/S_0}\otimes_{\mathscr O_{S_0}}\mathscr O_{Y_0},
     \mathscr J\mathscr O_{E\times_S Y}
   \right)\\
 &=\operatorname{Hom}_{\mathscr O_{E\times_S Y}}
   \left(
     \Omega^1_{E/S}\otimes_{\mathscr O_S}\mathscr O_Y,
     \mathscr J\mathscr O_{E\times_S Y}
   \right)\\
 &=\Gamma\!\left(
     E\times_S Y,
     \widetilde{\operatorname{Hom}}_{\mathscr O_{E\times_S Y}}
     \left(
       \Omega^1_{E/S}\otimes_{\mathscr O_S}\mathscr O_Y,
       \mathscr J\mathscr O_{E\times_S Y}
     \right)
   \right).
\end{align*}

Suppose that \(Y\) is flat over \(S\). There are then isomorphisms
\[
 \mathscr J\mathscr O_{E\times_S Y}
 \xleftarrow{\ \sim\ }
 (\mathscr J\mathscr O_E)\otimes_{\mathscr O_S}\mathscr O_Y
 \xrightarrow{\ \sim\ }
 (\mathscr J\mathscr O_E)\otimes_{\mathscr O_{S_0}}\mathscr O_{Y_0}.
\]
Suppose in addition that \(E\) is of finite presentation over \(S\).
Then \(\Omega^1_{E/S}\) is an \(\mathscr O_E\)-module of finite
presentation (EGA~IV\textsubscript{4}, 16.4.22), and therefore, by
EGA~0\textsubscript{I}, 6.7.6,
\begin{align*}
 &\widetilde{\operatorname{Hom}}_{\mathscr O_{E\times_S Y}}
 \left(
   \Omega^1_{E/S}\otimes_{\mathscr O_S}\mathscr O_Y,
   (\mathscr J\mathscr O_E)\otimes_{\mathscr O_S}\mathscr O_Y
 \right)\\
 &\qquad\simeq
 \widetilde{\operatorname{Hom}}_{\mathscr O_E}
 \left(\Omega^1_{E/S},\mathscr J\mathscr O_E\right)
 \otimes_{\mathscr O_S}\mathscr O_Y.
\end{align*}
Let \(\pi:E\to S\) and \(g:Y\to S\) be the structural morphisms.
Apply \SGARefLink{sga3:iii:paragraph:0-13}{0.13} to the diagram

\SGAThreeNativeDiagram[0.42\textwidth]
  {native_diagrams/exp3/diagram_027_exp3_p021_remark014_structural_base_change_square.tex}
  {The structural base-change square used in \SGARefLink{sga3:iii:remark:0-14}{Remark 0.14}.}
  {fig:exp3-remark-014-structural-base-change}
  {Polo--Gille re-edition, Expos\'e III}

and to the \(\mathscr O_E\)-module
\[
 \mathscr F=
 \widetilde{\operatorname{Hom}}_{\mathscr O_E}
 \left(\Omega^1_{E/S},\mathscr J\mathscr O_E\right).
\]
We obtain
\begin{align*}
 \Gamma(E\times_S Y,g'^*\mathscr F)
 &=\Gamma(Y,\pi'_*g'^*\mathscr F)\\
 &=\Gamma(Y,g^*\pi_*\mathscr F)\\
 &=W(\pi_*\mathscr F)(Y).
\end{align*}
Thus, if \(E\) is of finite presentation over \(S\), we have
\begin{equation}
 L'_X=
 W\!\left(
   \pi_*\widetilde{\operatorname{Hom}}_{\mathscr O_E}
   (\Omega^1_{E/S},\mathscr J\mathscr O_E)
 \right)
 \tag{0.14.1}\label{sga3:iii:remark:0-14:equation:0-14-1}
\end{equation}
on the category of \(S\)-schemes flat over \(S\). Notice also that the
module to which \(W\) is applied is quasi-coherent, by EGA~I, 9.1.1
and 9.2.1.

Let \(L_0\) be the \(S_0\)-functor\footnote{Editor's note: What follows
has been added.}
\[
 L_0=
 W\!\left(
   \pi_{0*}\widetilde{\operatorname{Hom}}_{\mathscr O_{E_0}}
   (\Omega^1_{E_0/S_0},\mathscr J\mathscr O_E)
 \right).
\]
Returning to the definition of \(L'_X(Y)\), and using the isomorphism
\[
 \mathscr J\mathscr O_{E\times_S Y}
 \simeq(\mathscr J\mathscr O_E)
       \otimes_{\mathscr O_{S_0}}\mathscr O_{Y_0},
\]
the same reasoning gives
\[
 L'_X(Y)=L_0(Y_0)=L_0(Y\times_S S_0)
 =\left(\prod_{S_0/S}L_0\right)(Y).
\]
Consequently, on the category of \(S\)-schemes flat over \(S\),
\[
 L'_X=
 \prod_{S_0/S}
 W\!\left(
   \pi_{0*}\widetilde{\operatorname{Hom}}_{\mathscr O_{E_0}}
   (\Omega^1_{E_0/S_0},\mathscr J\mathscr O_E)
 \right).
\]
It is not obvious that the action of \(X\) on \(L'_X\) defined in \SGARefLink{sga3:iii:corollary:0-12}{0.12}
comes from an action of
\(X_0=\underline{\operatorname{Aut}}_{S_0}(E_0)\) on \(L_0\).
This is nevertheless the case when \(E\) is also flat over \(S\).

Indeed, in that case there are canonical isomorphisms
\[
 \mathscr J\mathscr O_E
 \simeq\mathscr J\otimes_{\mathscr O_S}\mathscr O_E
 \simeq\mathscr J\otimes_{\mathscr O_{S_0}}\mathscr O_{E_0},
\]
and hence
\[
 L_0\simeq
 W\!\left(
   \pi_{0*}\widetilde{\operatorname{Hom}}_{\mathscr O_{E_0}}
   \left(
     \Omega^1_{E_0/S_0},
     \mathscr J\otimes_{\mathscr O_{S_0}}\mathscr O_{E_0}
   \right)
 \right),
\]
and
\begin{equation}
 L'_X=
 \prod_{S_0/S}
 W\!\left(
   \pi_{0*}\widetilde{\operatorname{Hom}}_{\mathscr O_{E_0}}
   \left(
     \Omega^1_{E_0/S_0},
     \mathscr J\otimes_{\mathscr O_{S_0}}\mathscr O_{E_0}
   \right)
 \right).
 \tag{0.14.2}\label{sga3:iii:remark:0-14:equation:0-14-2}
\end{equation}
For every \(S_0\)-scheme \(T\), we then have
\[
 L_0(T)\simeq
 \operatorname{Hom}_{\mathscr O_{E_0\times_{S_0}T}}
 \left(
   \Omega^1_{E_0\times_{S_0}T/T},
   \mathscr J\otimes_{\mathscr O_{S_0}}
     \mathscr O_{E_0\times_{S_0}T}
 \right).
\]
Moreover,
\[
 \operatorname{Hom}_{S_0}(T,X_0)
 =\operatorname{Aut}_T(E_0\times_{S_0}T)
\]
acts on \(L_0(T)\) by transport of structure, functorially in \(T\).
Finally, for every \(S\)-scheme \(Y\) flat over \(S\), the action by
transport of structure of
\[
 \operatorname{Hom}_S(Y,X)=\operatorname{Aut}_Y(E\times_S Y)
\]
on \(L'_X(Y)=L_0(Y_0)\) factors through
\(\operatorname{Aut}_{Y_0}(E_0\times_{S_0}Y_0)\).

Finally, let us extract the following two propositions from SGA~1,
Expos\'e~III.

\medskip
\noindent\SGAInlineRefTarget{sga3:iii:proposition:0-15}\textbf{Proposition 0.15.}
\textit{(SGA~1, Expos\'e~III, 6.8)}
\textit{For every affine \(S_{\mathscr J}\)-scheme \(Y\) smooth over
\(S_{\mathscr J}\), there exists an \(S\)-scheme \(X\), smooth over
\(S\), such that \(X\times_S S_{\mathscr J}\simeq Y\); such an \(X\) is
unique up to a non-unique isomorphism.}

\medskip
\noindent\SGAInlineRefTarget{sga3:iii:proposition:0-16}\textbf{Proposition 0.16.}
\textit{(SGA~1, Expos\'e~III, 5.5)}
\textit{Let \(X\) be an \(S\)-scheme smooth over \(S\). For every affine
\(S\)-scheme \(Y\), the canonical map}
\[
 p_X(Y):\operatorname{Hom}_S(Y,X)
 \longrightarrow\operatorname{Hom}_S(Y,X^+)
 =\operatorname{Hom}_{S_{\mathscr J}}(Y_{\mathscr J},X_{\mathscr J})
\]
\textit{is surjective.}

\medskip
\noindent\SGAInlineRefTarget{sga3:iii:corollary:0-17}\textbf{Corollary 0.17.}
\textit{Let \(E\) be an affine \(S\)-scheme smooth over \(S\), and put
\(X=\underline{\operatorname{Aut}}_S(E)\). For every affine \(S\)-scheme
\(Y\), the canonical map}
\begin{align*}
 \operatorname{Aut}_Y(E\times_S Y)
  =\operatorname{Hom}_S(Y,X)
 &\longrightarrow \operatorname{Hom}_S(Y,X^+)\\
  &=\operatorname{Aut}_{Y_{\mathscr J}}
    (E_{\mathscr J}\times_{S_{\mathscr J}}Y_{\mathscr J})
\end{align*}
\textit{is surjective.}

Indeed, \(Y\times_S E\) is affine over the affine scheme \(Y\), hence
is affine. Applying \SGARefLink{sga3:iii:proposition:0-16}{0.16}, we see that every \(S_{\mathscr J}\)-morphism
\[
 Y_{\mathscr J}\times_{S_{\mathscr J}}E_{\mathscr J}
 \longrightarrow E_{\mathscr J}
\]
extends to an \(S\)-morphism
\[
 Y\times_S E\longrightarrow E.
\]
In other words, every
\(Y_{\mathscr J}\)-endomorphism of
\(Y_{\mathscr J}\times_{S_{\mathscr J}}E_{\mathscr J}\) lifts to a
\(Y\)-endomorphism of \(Y\times_S E\). Then \SGARefLink{sga3:iii:corollary:0-7}{0.7}~\SGARefLink{sga3:iii:corollary:0-7:item:a}{(a)} shows that every
\(Y_{\mathscr J}\)-automorphism of
\(Y_{\mathscr J}\times_{S_{\mathscr J}}E_{\mathscr J}\) lifts to a
\(Y\)-automorphism of \(Y\times_S E\), which is the asserted property.

